\documentclass[10pt,a4paper]{article}

\usepackage[a4paper,top=17mm,bottom=18mm,left=16mm,right=16mm,headheight=14pt]{geometry}
\usepackage{fontspec}
\usepackage{xcolor}
\usepackage{array}
\usepackage{tabularx}
\usepackage{enumitem}
\usepackage{fancyhdr}
\usepackage{hyperref}
\usepackage{needspace}

\setmainfont{Arial}
\setsansfont{Arial}

\definecolor{Primary}{HTML}{235347}
\definecolor{Accent}{HTML}{E98B2A}
\definecolor{Soft}{HTML}{EEF4F1}
\definecolor{LineGray}{HTML}{B7C3BE}
\definecolor{TextGray}{HTML}{4E5A56}

\newcommand{\documentlabel}{Pedoman Wawancara Terbuka Karyawan Nasi Gerilya}

\hypersetup{
  colorlinks=true,
  linkcolor=Primary,
  urlcolor=Primary,
  pdftitle={Pedoman Wawancara Terbuka Karyawan Nasi Gerilya},
  pdfauthor={Instrumen Penelitian Skripsi}
}

\pagestyle{fancy}
\fancyhf{}
\fancyhead[L]{\small\textcolor{Primary}{\documentlabel}}
\fancyhead[R]{\small\textcolor{TextGray}{Versi Juli 2026}}
\fancyfoot[C]{\small\textcolor{TextGray}{Halaman \thepage}}
\renewcommand{\headrulewidth}{0.4pt}
\renewcommand{\footrulewidth}{0pt}

\setlength{\parindent}{0pt}
\setlength{\parskip}{4pt}
\renewcommand{\arraystretch}{1.2}
\hyphenpenalty=10000
\exhyphenpenalty=10000
\emergencystretch=3em
\sloppy
\newcolumntype{Y}{>{\raggedright\arraybackslash}X}

\newcommand{\sectionbar}[1]{%
  \Needspace{4\baselineskip}
  \vspace{4mm}
  \noindent\colorbox{Soft}{%
    \parbox{\dimexpr\linewidth-2\fboxsep\relax}{%
      \large\bfseries\textcolor{Primary}{#1}
    }%
  }\par\vspace{2mm}
}

\newcommand{\subsectionbar}[1]{%
  \Needspace{3\baselineskip}
  \vspace{2.5mm}
  {\normalsize\bfseries\textcolor{Accent}{#1}}\par
  \textcolor{Accent}{\rule{0.28\linewidth}{0.7pt}}\par\vspace{1mm}
}

\newcommand{\answerline}{%
  \par\vspace{0.9mm}
  \noindent\textcolor{LineGray}{\rule{\linewidth}{0.35pt}}\par\vspace{3mm}
}

\newcommand{\question}[1]{%
  \Needspace{4.5\baselineskip}
  \item #1
  \answerline
}

\newcommand{\widefield}[1]{%
  \Needspace{3\baselineskip}
  \textbf{#1}\par
  \noindent\textcolor{LineGray}{\rule{\linewidth}{0.35pt}}\par\vspace{2.7mm}
}

\begin{document}

\begin{titlepage}
\thispagestyle{empty}
\vspace*{10mm}
{\Huge\bfseries\textcolor{Primary}{Pedoman Wawancara}}\par
{\Huge\bfseries\textcolor{Primary}{Terbuka Karyawan}}\par
{\Huge\bfseries\textcolor{Primary}{Nasi Gerilya}}\par
\vspace{4mm}
{\Large\textcolor{Accent}{Instrumen Angket Terbuka Karyawan}}\par
\vspace{4mm}
{\large\textcolor{TextGray}{Versi terbaru: 01 Juli 2026}}\par
\vspace{7mm}
\textcolor{Accent}{\rule{0.92\linewidth}{1.4pt}}\par
\vspace{10mm}

\begin{tabularx}{\linewidth}{>{\bfseries\color{Primary}}p{36mm} Y}
Topik penelitian & Strategi pemasaran Rumah Makan Nasi Gerilya pada platform GrabFood untuk meningkatkan penjualan berdasarkan analisis SWOT.\\
Sasaran informan & Karyawan dapur, kasir, pelayanan, atau karyawan yang terlibat dalam operasional Nasi Gerilya dan pesanan GrabFood.\\
Bentuk instrumen & Pertanyaan terbuka. Jawaban dapat ditulis langsung oleh informan atau dicatat oleh peneliti saat wawancara.\\
Kerahasiaan & Identitas informan dijaga kerahasiaannya dan data digunakan hanya untuk keperluan penelitian.\\
\end{tabularx}

\vspace{9mm}
\colorbox{Soft}{%
  \parbox{\dimexpr\linewidth-2\fboxsep\relax}{%
    \raggedright
    \textbf{\textcolor{Primary}{Petunjuk penggunaan:}}\par
    Untuk karyawan dapur gunakan bagian A + B + C + E.\par
    Untuk karyawan kasir gunakan bagian A + B + D + E.
  }%
}

\vfill
{\small\textcolor{TextGray}{Jawablah berdasarkan pengalaman kerja sehari-hari di Nasi Gerilya. Tidak ada jawaban benar atau salah. Ceritakan contoh nyata jika memungkinkan.}}
\vspace*{8mm}
\end{titlepage}

\sectionbar{Identitas Singkat Informan}
\widefield{Kode informan}
\widefield{Tanggal wawancara atau pengisian}
\widefield{Bagian kerja}
\widefield{Shift atau jam kerja}
\widefield{Catatan peneliti}

\sectionbar{A. Pertanyaan Pembuka untuk Kedua Karyawan}
\begin{enumerate}[label=\arabic*.,leftmargin=*,itemsep=1.5mm]
\question{Sejak kapan Anda bekerja di Nasi Gerilya?}
\question{Bagaimana awalnya Anda bisa bekerja di Nasi Gerilya?}
\question{Apa tugas utama Anda sehari-hari?}
\question{Biasanya Anda mulai bekerja dari jam berapa sampai jam berapa?}
\question{Kalau baru mulai buka, apa saja yang pertama kali dikerjakan?}
\question{Kalau mau tutup, apa saja yang biasanya dibereskan?}
\question{Bagaimana pembagian kerja antara bagian dapur dan bagian kasir?}
\question{Menurut Anda, bagian pekerjaan apa yang paling berat di Nasi Gerilya? Ceritakan alasannya.}
\end{enumerate}

\sectionbar{B. Alur Kerja dan Kondisi Pesanan}
\begin{enumerate}[label=\arabic*.,leftmargin=*,itemsep=1.5mm]
\question{Ceritakan alur kerja dari pesanan masuk sampai makanan diberikan ke pelanggan atau driver.}
\question{Kapan biasanya pesanan paling ramai?}
\question{Kapan biasanya pesanan paling sepi?}
\question{Kalau pesanan sedang ramai, apa yang biasanya terjadi di tempat kerja?}
\question{Kalau pesanan sedang ramai, bagaimana Anda dan karyawan lain membagi pekerjaan?}
\question{Pernahkah pesanan menumpuk? Ceritakan bagaimana kejadiannya.}
\question{Kesalahan apa yang pernah terjadi saat pesanan sedang banyak?}
\question{Kalau ada kesalahan pesanan, biasanya siapa yang mengetahui lebih dulu?}
\question{Bagaimana cara memperbaiki kesalahan pesanan tersebut?}
\question{Menurut Anda, apa yang perlu diperbaiki supaya pekerjaan lebih lancar?}
\end{enumerate}

\sectionbar{C. Pertanyaan Khusus Karyawan Bagian Dapur}
\subsectionbar{Persiapan dan Memasak}
\begin{enumerate}[label=\arabic*.,leftmargin=*,itemsep=1.5mm]
\question{Apa saja bahan makanan yang biasanya disiapkan sebelum mulai jualan?}
\question{Bagaimana cara menentukan bahan yang harus dimasak lebih dulu?}
\question{Menu apa yang paling sering dibuat?}
\question{Menu apa yang paling lama dibuat?}
\question{Menu apa yang paling mudah dibuat?}
\question{Saat membuat pesanan, bagaimana Anda tahu pesanan mana yang harus dikerjakan dulu?}
\question{Kalau ada banyak pesanan sekaligus, bagaimana cara Anda mengatur urutannya?}
\question{Pernahkah makanan harus dibuat ulang? Ceritakan penyebabnya.}
\question{Apa yang biasanya membuat proses memasak menjadi lambat?}
\question{Apa yang membantu pekerjaan dapur menjadi lebih cepat?}
\end{enumerate}

\subsectionbar{Rasa, Porsi, dan Kualitas Makanan}
\begin{enumerate}[label=\arabic*.,leftmargin=*,itemsep=1.5mm]
\question{Bagaimana cara Anda menjaga rasa makanan supaya tetap sama?}
\question{Bagaimana cara Anda menentukan porsi untuk setiap pesanan?}
\question{Apa yang Anda lakukan supaya porsi tidak berbeda-beda?}
\question{Bagaimana cara memastikan makanan sudah matang dan siap diberikan?}
\question{Apa yang biasanya dicek sebelum makanan dibungkus?}
\question{Kalau ada pelanggan minta catatan khusus, bagaimana dapur menanganinya?}
\question{Apa catatan khusus dari pelanggan yang paling sering diterima?}
\question{Menurut Anda, bagian mana dari makanan Nasi Gerilya yang paling disukai pelanggan?}
\question{Apa keluhan pelanggan tentang makanan yang pernah Anda dengar?}
\question{Kalau makanan dikirim lewat GrabFood, apa yang perlu diperhatikan agar tetap enak sampai ke pelanggan?}
\end{enumerate}

\subsectionbar{Bahan Baku dan Stok}
\begin{enumerate}[label=\arabic*.,leftmargin=*,itemsep=1.5mm]
\question{Bahan apa yang paling cepat habis?}
\question{Bagaimana Anda tahu kalau stok bahan mulai menipis?}
\question{Kalau bahan hampir habis tapi masih ada pesanan, apa yang biasanya dilakukan?}
\question{Pernahkah ada menu yang tidak bisa dijual karena bahan habis? Ceritakan.}
\question{Bagaimana cara menyimpan bahan supaya tetap bagus?}
\question{Bahan apa yang paling sulit dijaga kualitasnya?}
\question{Kalau harga bahan naik, apa yang biasanya terasa di pekerjaan dapur?}
\question{Peralatan dapur apa yang paling sering digunakan?}
\question{Peralatan apa yang menurut Anda masih kurang atau perlu diganti?}
\end{enumerate}

\subsectionbar{Kebersihan Dapur}
\begin{enumerate}[label=\arabic*.,leftmargin=*,itemsep=1.5mm]
\question{Bagaimana cara menjaga kebersihan dapur saat bekerja?}
\question{Bagaimana cara membersihkan alat masak setelah dipakai?}
\question{Kalau sedang ramai, bagian kebersihan apa yang paling sulit dijaga?}
\question{Apa yang biasanya dilakukan supaya makanan tetap bersih saat dimasak dan dibungkus?}
\end{enumerate}

\subsectionbar{Kerja Sama dengan Kasir}
\begin{enumerate}[label=\arabic*.,leftmargin=*,itemsep=1.5mm]
\question{Bagaimana kasir menyampaikan pesanan ke dapur?}
\question{Informasi apa yang paling penting disampaikan kasir ke dapur?}
\question{Pernahkah pesanan kurang jelas dari kasir atau aplikasi? Ceritakan.}
\question{Kalau makanan sudah selesai, bagaimana dapur memberi tahu kasir?}
\question{Menurut Anda, apa yang perlu diperbaiki dari kerja sama dapur dan kasir?}
\end{enumerate}

\sectionbar{D. Pertanyaan Khusus Karyawan Bagian Kasir/Pelayanan/GrabFood}
\subsectionbar{Tugas Kasir dan Pelayanan}
\begin{enumerate}[label=\arabic*.,leftmargin=*,itemsep=1.5mm]
\question{Apa saja tugas Anda sebagai kasir?}
\question{Bagaimana cara Anda menerima pesanan dari pelanggan yang datang langsung?}
\question{Bagaimana cara Anda mencatat pesanan pelanggan?}
\question{Bagaimana cara Anda menyampaikan pesanan ke dapur?}
\question{Apa yang biasanya ditanyakan pelanggan sebelum membeli?}
\question{Bagaimana cara Anda menjawab pelanggan yang bingung memilih menu?}
\question{Menu apa yang paling sering ditanyakan pelanggan?}
\question{Menu apa yang paling sering dibeli pelanggan langsung?}
\question{Apa keluhan pelanggan yang pernah Anda terima?}
\question{Kalau pelanggan komplain, apa yang biasanya Anda lakukan?}
\end{enumerate}

\subsectionbar{Pesanan GrabFood}
\begin{enumerate}[label=\arabic*.,leftmargin=*,itemsep=1.5mm]
\question{Ceritakan bagaimana pesanan GrabFood masuk ke Nasi Gerilya.}
\question{Apa yang Anda lakukan setelah pesanan GrabFood masuk?}
\question{Bagaimana cara mengecek apakah menu yang dipesan masih tersedia?}
\question{Bagaimana cara menyampaikan pesanan GrabFood ke dapur?}
\question{Bagaimana cara memastikan catatan khusus dari pelanggan tidak terlewat?}
\question{Apa catatan khusus pelanggan GrabFood yang paling sering muncul?}
\question{Bagaimana cara mengecek makanan sebelum diberikan ke driver?}
\question{Bagaimana cara memastikan pesanan tidak tertukar?}
\question{Pernahkah pesanan GrabFood salah atau kurang? Ceritakan kejadiannya.}
\question{Kalau pesanan GrabFood dibatalkan, apa yang biasanya dilakukan?}
\question{Kalau aplikasi GrabFood bermasalah, apa yang biasanya terjadi?}
\question{Apa kesulitan paling sering saat mengurus pesanan GrabFood?}
\end{enumerate}

\subsectionbar{Driver GrabFood}
\begin{enumerate}[label=\arabic*.,leftmargin=*,itemsep=1.5mm]
\question{Bagaimana biasanya komunikasi dengan driver GrabFood?}
\question{Kalau driver datang sebelum makanan selesai, apa yang Anda lakukan?}
\question{Kalau makanan sudah selesai tapi driver belum datang, apa yang Anda lakukan?}
\question{Pernahkah ada masalah dengan driver? Ceritakan.}
\question{Apa yang membuat proses penyerahan pesanan ke driver berjalan lancar?}
\end{enumerate}

\subsectionbar{Menu, Harga, dan Tampilan di GrabFood}
\begin{enumerate}[label=\arabic*.,leftmargin=*,itemsep=1.5mm]
\question{Menu apa yang paling sering dipesan lewat GrabFood?}
\question{Menurut pengamatan Anda, kenapa menu itu sering dipesan?}
\question{Menu apa yang jarang dipesan lewat GrabFood?}
\question{Menurut Anda, kenapa menu itu jarang dipesan?}
\question{Apa yang pernah dikatakan pelanggan tentang harga makanan?}
\question{Apa yang pernah dikatakan pelanggan tentang porsi makanan?}
\question{Apa yang pernah dikatakan pelanggan tentang kemasan?}
\question{Menurut Anda, apakah foto menu di GrabFood sudah sesuai dengan makanan aslinya? Ceritakan.}
\question{Informasi apa yang menurut Anda perlu ditambahkan di GrabFood supaya pelanggan lebih mudah memilih?}
\end{enumerate}

\subsectionbar{Promosi dan Pembelian}
\begin{enumerate}[label=\arabic*.,leftmargin=*,itemsep=1.5mm]
\question{Kalau sedang ada promo, apa yang biasanya terjadi pada pesanan?}
\question{Promo seperti apa yang paling sering membuat pelanggan membeli?}
\question{Dari mana biasanya pelanggan tahu tentang Nasi Gerilya?}
\question{Pernahkah pelanggan datang atau memesan karena melihat promo? Ceritakan.}
\question{Menurut Anda, apa yang membuat pelanggan mau membeli lagi?}
\question{Menurut Anda, apa yang bisa membuat pelanggan pindah membeli ke tempat lain?}
\end{enumerate}

\sectionbar{E. Pertanyaan Penutup untuk Kedua Karyawan}
\begin{enumerate}[label=\arabic*.,leftmargin=*,itemsep=1.5mm]
\question{Menurut Anda, apa kelebihan utama Nasi Gerilya?}
\question{Menurut Anda, apa kekurangan yang paling perlu diperbaiki?}
\question{Apa saran Anda supaya pekerjaan karyawan lebih mudah?}
\question{Apa saran Anda supaya pelanggan lebih puas?}
\question{Apa saran Anda supaya pesanan GrabFood lebih lancar?}
\question{Apa harapan Anda untuk Nasi Gerilya ke depannya?}
\question{Apakah ada hal penting lain yang belum ditanyakan tetapi ingin Anda ceritakan?}
\end{enumerate}

\end{document}
